\subsection{Excluding the remaining class-VII minimal models}

\begin{statement}{P4.37}{Proved}
No irreducible singular support in Part~4.2 has normalization
whose minimal resolution has a class-VII minimal model with
positive second Betti number. Consequently every class-VII
normalization occurring here is a smooth minimal Hopf surface
and has Euler number zero.
\end{statement}
\begin{proof}
The numerically anticanonical classification used in P4.23
and P4.26 reduces this question to an intermediate Kato
minimal model $T_0$: the Inoue--Hirzebruch and parabolic
alternatives were excluded there by their reduced
normal-crossing conductor and the Chern identity.
Write $n=b_2(T_0)>0$, let $A_0$ be its maximal reduced
curve divisor, and let $D_0$ be the effective integral
pole divisor pushed down from the normalization resolution.
Thus $K_{T_0}+D_0$ is numerically zero. The support of
$D_0$ is $A_0$.

We state the precise additional primary input.
Favre, \emph{Dynamique des applications rationnelles},
Theorem~1.2.24(1)--(2), printed p.~35, says that the
$n$ rational curve classes span $H^2(T_0;\mathbb Q)$
and that their intersection form is negative definite.
Theorem~1.2.30, printed p.~42, and its proof on
pp.~44 and~46 give a holomorphic foliation with conormal
line
\[
 N_{\mathcal F}^*=\mathcal O_{T_0}(-A_0)
                       \otimes\text{a flat line}.
\]
Theorem~1.2.32 and its proof on p.~46 show that every
component of $A_0$ is invariant and that its singularities
are precisely the intersection points, with local
first integrals $z^\alpha w^\beta$, $\alpha,\beta>0$.
The invariant-curve and order-one assertions are also
Dloussky--Oeljeklaus, Theorem~3.1, pp.~1528--1529.
The actual source URLs and exact theorem locators are
recorded in the bibliography and literature log.

It follows from the determinant of tangent and normal
lines that the foliation tangent line has numerical
first Chern class $-K_{T_0}-A_0$. The local generator
at an intersection is
$\beta z\partial_z-\alpha w\partial_w$.
Blowing up a regular point or one of these saddles
again gives only regular points and saddles with
opposite-sign nonzero eigenvalues. Thus every actual
blowup centre has the same two possibilities
$\varepsilon=0,1$ as in P4.36, with the same formulas.

Let $T\to T_0$ contain $k$ blowups. Write
$c_1(L)=-K_T+V$ for the saturated foliation line;
initially $V_0=-A_0$, and
$V'=\mu^*V+\varepsilon E$ at a blowup.
The reduced divisor $A_0$ is a rational cycle with
rational trees attached. Its normalization has
$n$ rational components and $n$ nodes, including
self-nodes, so $\chi(\mathcal O_{A_0})=0$.
Adjunction, or the divisor exact sequence on $T_0$,
therefore gives
\[
 (K_{T_0}+A_0)A_0=0.
\]
Hence $K_{T_0}V_0-V_0^2=0$.
Under each blowup this quantity changes by
$\varepsilon^2-\varepsilon=0$.
Since $K_{T_0}^2=-n$, we obtain
\[
 K_TV-V^2=0,\qquad K_T^2=-n-k.
\]

Define $B=\rho^*TY$, the divisorial determinant zero
$R$, the divisorial foliation zero $Z$, and the
nonnegative lengths $z_0,z_1$ exactly as in P4.36.
Its saturated sheaf calculation, including the
Hartogs extension of the rank-one double dual,
now gives
\[
 \mathcal B(R,Z)
 =R^2-K_TR-VR-ZR-K_TZ+2VZ+Z^2
 =n+k+z_0+z_1.                                  \tag{P4.37.1}
\]
Every reversed blowup changes $\mathcal B$ by at
most one, by the unrestricted integer identity
(P4.36.2). It remains to evaluate the initial term.

Push $R,Z$ down to effective divisors $R_0,Z_0$
on $T_0$. Their supports are contained in $A_0$:
the derivative is generically injective away from
the resolved conductor, and in any event all
compact curves of the intermediate surface are
components of its maximal divisor. Put $D=D_0$,
$A=A_0$. Using $K=-D$ numerically and $V=-A$ gives
\[
 \mathcal B_0
 =R_0^2+DR_0+AR_0-Z_0R_0+DZ_0-2AZ_0+Z_0^2.
\]
For divisor classes define
$q(X,Y)=X^2-XY+Y^2$.
Adjunction above gives $DA=A^2$, while $D^2=-n$.
Expansion consequently gives the exact completion
\[
 \mathcal B_0
       =n+q(R_0+D,Z_0+D-A).                     \tag{P4.37.2}
\]
The form on divisor classes is negative definite.
Moreover
\[
 q(X,Y)=\tfrac12\bigl(X^2+Y^2+(X-Y)^2\bigr),
\]
so it is nonpositive, and it vanishes only for
$X=Y=0$. The first argument $R_0+D$ is nonzero:
it is a nonzero effective combination of the
$n$ independent curve classes. Thus
$\mathcal B_0<n$.
After the $k$ blowups we still have
$\mathcal B(R,Z)<n+k$, contradicting (P4.37.1).
This excludes the intermediate minimal model.

All class-VII cases with positive second Betti
number have now been excluded. The zero-Betti
Inoue case has no conductor curves and was
excluded in P4.26. P4.36 treats the remaining
Hopf case and proves the final assertion.
\end{proof}
